% =====================================================================
% Supplementary Material -- sections moved out of the main text on the
% supervisor's instruction (2026-09-26): run-to-run repeats and practical
% identifiability. Text is the authors' own from the main-text draft
% (PDF version 6, Sections 2.4.2 last paragraph, 2.4.4, 3.4, 3.5 and
% Tables 5-6), with only the cross-references changed to fit the
% supplement. Paste into the supplement file on Overleaf.
% Requires in the preamble: booktabs, amsmath.
% =====================================================================

\section{Run-to-run repeats of the calibration pipeline}
\label{sec:S-repeats}

\subsection{Method}

The complete pipeline was executed three times with the multi-start seeds
101, 202 and 303, fixed in advance, with clinical inputs and all other
settings held constant. The seed enters only through the scrambling of
starting points 2--6; starting point 1 is deterministic. All three runs are
reported.

\subsection{Result}

The three runs returned the same accepted parameter set, so the pipeline is
exactly reproducible under the seeds tested. The reason is visible in the
multi-start record (Table~\ref{tab:S-multistart}): the clinically seeded
starting point is deterministic and is not affected by the seed, it was the
only start with no hard-target failure and the lowest governed error in runs
101 and 202, and it won the lexicographic selection in all three runs. The
seed changes only the five scrambled starting points.

Variability therefore exists between starting points rather than between
runs. Across the 18 starts, the governed RMSE ranged from 0.0909 to 0.1509.
One scrambled start in run 303 reached 0.0909, below the accepted value of
0.1047, but failed one hard target and was rejected by the selection rule.
This is the intended behaviour of the rule, and it means the reported result
is not the lowest error attained, but the lowest error attained without a
hard-target failure.

Because the accepted candidate does not vary, the three repeats do not
quantify uncertainty. The measurement-noise simulation of the main text is
the analysis intended for that purpose.

\begin{table}[h]
\centering
\caption{Governed RMSE of the six starting points in each run, with the
number of hard-target failures and gate failures in parentheses. The
selected start is in bold.}
\label{tab:S-multistart}
\begin{tabular}{lccc}
\toprule
Start & Run 101 & Run 202 & Run 303 \\
\midrule
Clinically seeded & \textbf{0.1047 (0, 2)} & \textbf{0.1047 (0, 2)} & \textbf{0.1047 (0, 2)} \\
Sobol' 1 & 0.1509 (1, 5) & 0.1509 (1, 5) & 0.1509 (1, 5) \\
Sobol' 2 & 0.1165 (1, 4) & 0.1288 (1, 3) & 0.1209 (1, 3) \\
Sobol' 3 & 0.1211 (1, 3) & 0.1184 (1, 2) & 0.1509 (1, 5) \\
Sobol' 4 & 0.1345 (0, 2) & 0.1392 (0, 2) & 0.1340 (1, 4) \\
Sobol' 5 & 0.1229 (0, 2) & 0.1509 (1, 5) & 0.0909 (1, 2) \\
\bottomrule
\end{tabular}
\end{table}

\section{Practical identifiability}
\label{sec:S-identifiability}

\subsection{Method}

Practical identifiability was assessed at the calibrated operating point from
the scaled sensitivity matrix \cite{pope2009,olufsen2013}
\begin{equation}
S_{ij} = \frac{\partial y_i}{\partial \theta_j}\,\frac{\theta_j}{\sigma_i},
\end{equation}
with $\sigma_i$ the declared measurement uncertainty of target $i$.
Parameters with small column norms are weakly constrained; pairs with
$|\rho| \ge 0.9$ between columns were flagged as collinear. The condition
number $\mathrm{cond}(S)$ was evaluated for every subset of size $k \le p$
and the best-conditioned subset at each size reported as a candidate reduced
set \cite{vanosta2020}. Subsets were scored by linear algebra alone, without
re-simulation or recalibration.

\subsection{Result}

With 9 governed targets and 7 fitted parameters the pre-closure problem has
2 degrees of freedom. At the calibrated operating point, $\mathrm{cond}(S)$
of the full seven-parameter set was 61.9, identical in all three runs. No
parameter pair was collinear at the $|\rho| \ge 0.9$ criterion; the strongest
correlation was $-0.73$ between the pulmonary resistance scale and
$C_{\mathrm{SAR}}$, followed by $+0.51$ between the systemic resistance scale
and $R_{\mathrm{SVEN}}$. The weakest column was $C_d$ (norm 3.72), followed by
$R_{\mathrm{SVEN}}$ (6.56); the pulmonary resistance scale was the
best-constrained (27.6).

Conditioning improved sharply on removal of one parameter
(Table~\ref{tab:S-cond}): the best six-parameter subset, which drops
$R_{\mathrm{SVEN}}$, had $\mathrm{cond}(S) = 9.68$. The seven-parameter set
retained here is therefore the least well-conditioned of the candidates, and
is kept because it is the set the screen selected, not because it is the best
conditioned. The reduced subsets were scored by linear algebra alone and were
not recalibrated; they are reported as candidates for a future run.

\begin{table}[h]
\centering
\caption{Best condition number of the scaled sensitivity matrix by subset
size ($N = 9$ governed targets). Values are identical in all three runs.}
\label{tab:S-cond}
\begin{tabular}{rrrl}
\toprule
$p$ & dof & $\mathrm{cond}(S)$ & Parameters \\
\midrule
7 & $+2$ & 61.90 & full retained set \\
6 & $+3$ & 9.68 & $R_{\mathrm{sys}}, R_{\mathrm{pul}}, C_{\mathrm{SAR}}, C_{\mathrm{PAR}}, E_{\mathrm{B,RV}}, C_d$ \\
5 & $+4$ & 4.00 & $R_{\mathrm{SVEN}}, C_{\mathrm{SAR}}, C_{\mathrm{PAR}}, E_{\mathrm{B,RV}}, C_d$ \\
4 & $+5$ & 2.65 & $R_{\mathrm{SVEN}}, C_{\mathrm{PAR}}, E_{\mathrm{B,RV}}, C_d$ \\
3 & $+6$ & 1.72 & $C_{\mathrm{SAR}}, C_{\mathrm{PAR}}, E_{\mathrm{B,RV}}$ \\
2 & $+7$ & 1.07 & $C_{\mathrm{PAR}}, E_{\mathrm{B,RV}}$ \\
\bottomrule
\end{tabular}
\end{table}
